 % Il progetto nasce dal template per il frontespizio realizzato da Marco Antonio Corallo, che ringrazio. Seguono alcuni commenti per evidenziare la presenza di alcuni pacchetti che mi sono stati utili per la stesura della tesi. Chiaramente, dipende tutto dal tipo di lavoro che uno vuole eseguire, che determina anche le diverse esigenze. Durante la stesura ho passato molto tempo su siti e forum a cercare di risolvere alcuni probelmi di formattazione, ma in generale Latex è stato piuttosto versatile. 

% Tipo di documento. L'uso di twoside implica che i capitoli inizino sempre con la prima pagina a sinistra, eventualmente lasciando una pagina vuota nel capitolo precedente. Se questa cosa è fastidiosa, è possibile rimuoverlo. 
\documentclass[a4paper, twoside,openright]{report}

% Dimensione dei margini
\usepackage[a4paper,top=3cm,bottom=3cm,left=3cm,right=3cm]{geometry}
% Dimensione del font
\usepackage[fontsize=13pt]{scrextend}
% Lingua del testo
\usepackage[provide=*,italian]{babel}
% Lingua per la bibliografia
\usepackage[fixlanguage]{babelbib}
% Codifica del testo
% Encoding del testo
\usepackage[T1]{fontenc}
% Permette di generare testo fittizio. Mi è stato utile
% per capire quale sarebbe stata l'impostazione del
% testo nella pagina prima che scrivessi un determinato paragrafo
\usepackage{lipsum}
% Per ruotare le immagini
\usepackage{rotating}
% Per modificare l'header delle pagine
\usepackage{fancyhdr}

% Librerie matematiche
\usepackage{amssymb}
\usepackage{amsmath}
\usepackage{amsthm}

% Uso delle immagini
\usepackage{graphicx}
% Uso dei colori
\usepackage[dvipsnames]{xcolor}
% Uso dei listing per il codice
\usepackage{listings}
% Per inserire gli hyperlinks tra i vari elementi del testo
\usepackage{hyperref}
% Diversi tipi di sottolineature
\usepackage[normalem]{ulem}

\usepackage{svg}


% -----------------------------------------------------------------

% Modifica lo stile dell'header
\pagestyle{fancy}
\fancyhf{}
\lhead{\rightmark}
\rhead{\textbf{\thepage}}
\fancyfoot{}
\setlength{\headheight}{12.5pt}

% Rimuove il numero di pagina all'inizio dei capitoli
\fancypagestyle{plain}{
  \fancyfoot{}
  \fancyhead{}
  \renewcommand{\headrulewidth}{0pt}
}

% Stile del codice
\lstdefinestyle{codeStyle}{
    % Colore dei commenti
    commentstyle=\color{teal},
    % Colore delle keyword
    keywordstyle=\color{Magenta},
    % Stile dei numeri di riga
    numberstyle=\tiny\color{gray},
    % Colore delle stringhe
    stringstyle=\color{violet},
    % Dimensione e stile del testo
    basicstyle=\ttfamily\footnotesize,
    % newline solo ai whitespaces
    breakatwhitespace=false,
    % newline si/no
    breaklines=true,
    % Posizione della caption, top/bottom
    captionpos=b,
    % Mantiene gli spazi nel codice, utile per l'indentazione
    keepspaces=true,
    % Dove visualizzare i numeri di linea
    numbers=left,
    % Distanza tra i numeri di linea
    numbersep=5pt,
    % Mostra gli spazi bianchi o meno
    showspaces=false,
    % Mostra gli spazi bianchi nelle stringhe
    showstringspaces=false,
    % Mostra i tab
    showtabs=false,
    % Dimensione dei tab
    tabsize=2
} \lstset{style=codeStyle}

% Stile di codice per dimensioni maggiori, in cui ho avuto bisogno di un testo più piccolo (ad esempio se si vuole inserire del codice che ha linee molto lunghe). Per usare questo stile piuttosto che il precedente, usare 

% \lstset{style=longBlock}
%  % inserire il codice...
% \lstset{style=codeStyle}

% Il secondo comando consente di tornare allo stile precedente
\lstdefinestyle{longBlock}{
    commentstyle=\color{teal},
    keywordstyle=\color{Magenta},
    numberstyle=\tiny\color{gray},
    stringstyle=\color{violet},
    basicstyle=\ttfamily\scriptsize,
    breakatwhitespace=false,
    breaklines=true,
    captionpos=b,
    keepspaces=true,
    numbers=left,
    numbersep=5pt,
    showspaces=false,
    showstringspaces=false,
    showtabs=false,
    tabsize=2
} \lstset{style=codeStyle}

% Togliendo il commento al comando che segue, si inseriscono nella bibliografia anche le fonti presenti in Bibliography.bib ma non citati direttamente con il comando \cite
% \nocite{*}

% Margini prima e dopo blocchi di codice, per avere più distanza
\lstset{aboveskip=20pt,belowskip=20pt}

% Modifica dello stile dei riferimenti, con il testo in cyano
\hypersetup{
    colorlinks,
    linkcolor=CornflowerBlue,
    citecolor=CornflowerBlue
}

% Aggiunti definizioni, teoremi, linea e listing
\newtheorem{definition}{Definizione}[section]
\newtheorem{theorem}{Teorema}[section]
\providecommand*\definitionautorefname{Definizione}
\providecommand*\theoremautorefname{Teorema}
\providecommand*{\listingautorefname}{Listing}
\providecommand*\lstnumberautorefname{Linea}

\raggedbottom

%\newcommand{\cgs}[1]{{\textcolor{brown}[\textcolor{red}{\bf{GS: }}{ \textcolor{brown}{#1]}}}}
%\newcommand{\cmc}[1]{{\textcolor{blue}[\textcolor{magenta}{\bf{MC: }}{ \textcolor{blue}{#1]}}}}

\colorlet{numb}{magenta!60!black}

\lstdefinelanguage{json}{
    numbers=left,
    string=[s]{"}{"},
    comment=[l]{:\ "},
    morecomment=[l]{:"},
    literate=
        *{0}{{{\color{numb}0}}}{1}
         {1}{{{\color{numb}1}}}{1}
         {2}{{{\color{numb}2}}}{1}
         {3}{{{\color{numb}3}}}{1}
         {4}{{{\color{numb}4}}}{1}
         {5}{{{\color{numb}5}}}{1}
         {6}{{{\color{numb}6}}}{1}
         {7}{{{\color{numb}7}}}{1}
         {8}{{{\color{numb}8}}}{1}
         {9}{{{\color{numb}9}}}{1}
}

% -----------------------------------------------------------------
\begin{document}

\begin{titlepage}
    \begin{figure}[!htb]
        \centering
        \includegraphics[keepaspectratio=true,scale=0.5]{chapters/Frontespizio/images/cherubinFrontespizio.eps}
    \end{figure}

    \begin{center}
        \LARGE{UNIVERSITÀ DI PISA}
        \vspace{5mm}
        \\ \large{DIPARTIMENTO DI INFORMATICA }
        \vspace{5mm}
        \\ \LARGE{Corso di Laurea Triennale in Informatica}
    \end{center}

    \vspace{15mm}
    \begin{center}
        {\LARGE{\bf Sviluppo e integrazione di dati relativi alla struttura di opere nel contesto di un sistema di gestione dei ponti\\ \vspace{5mm} Progetto di Tirocinio }}

    \end{center}
    \vspace{30mm}


    \begin{minipage}[t]{0.47\textwidth}
        {\large{Relatori:}{\normalsize\vspace{3mm}
                \bf\\ \large{Stefano Forti
                    \vspace{2mm}\\Paolo Tesser}}}
    \end{minipage}
    \hfill
    \begin{minipage}[t]{0.47\textwidth}\raggedleft
        {\large{Candidato:}{\normalsize\vspace{3mm} \bf\\ \large{ Octavian Alexandru\\ }}}
    \end{minipage}

    \vspace{20mm}
    \hrulefill
    \\\centering{\large{ANNO ACCADEMICO 2023/2024}}

\end{titlepage}
\cleardoublepage\phantomsection\pdfbookmark{Sommario}{Sommario}
\begingroup
\let\clearpage\relax
\let\cleardoublepage\relax
\let\cleardoublepage\relax

\chapter*{Sommario}

Il presente documento descrive il lavoro svolto durante il periodo di tirocinio formativo, della durata di circa trecento ore complessive, dal laureando Samuele Calugi presso l'azienda Geckosoft, sede di Pisa.
Durante il tirocinio era richiesto il raggiungimento di massima dei seguenti obiettivi:
\begin{itemize}
    \item studio autonomo degli Standard OGC (Web Map Service, Web Feature Service, Web Coverage Service, Web Map Tile Service) e dei formati dati geospaziali (GeoJSON, Shapefile, ecc.);
    \item sviluppo di un frontend interattivo in grado di mostrare mappe ottenute attraverso gli standard OGC e i formati dati geospaziali specificati;
    \item sviluppo di un backend per il fornimento delle mappe in formato conforme agli standard OGC e ai formati dati geospaziali richiesti;
    \item integrazione del backend e del frontend per creare un'applicazione completa e funzionale per la visualizzazione di dati geospaziali provenienti da diverse fonti.
    \item l'utilizzo di \textit{Angular}, per la parte front-end e \textit{.NET} per la parte back-end;
    \item documentazione dettagliata del lavoro svolto, incluse le tecnologie utilizzate, le scelte progettuali e le modalità di utilizzo del prodotto.
\end{itemize}

\endgroup

\vfill

\tableofcontents

% Rimuovere se non si vuole la tabella delle figure
\listoffigures

\include{chapters/1/chap1}
\include{chapters/2/chap2}
\include{chapters/3/chap3}
\include{chapters/4/chap4}

\appendix

% \chapter{Questa è un'appendice}

\lipsum[5]

\section*{Questa è un'altra sezione, ma non viene inserita nell'indice}

\lipsum[6]


\bibliographystyle{plain}
\bibliography{chapters/Bibliografia.bib}

\end{document}
% -----------------------------------------------------------------
